\documentclass[../../../include/open-logic-section]{subfiles}

\begin{document}

\olfileid{sth}{card-arithmetic}{fix}
\olsection{نقاط ثابتِ $\aleph$}

در \olref[spine][]{chap} گفتیم اصل جایگزینی به توجیه نیاز دارد، زیرا
سلسله‌مراتب را ناگزیر بسیار بلند می‌کند. اکنون که قدری حساب کاردینالی
آموخته‌ایم، می‌توانیم تصویری کوچک از بلندی این سلسله‌مراتب به دست دهیم.

آشکارا $0 < \aleph_0$ و $1 < \aleph_1$ و $2 < \aleph_2$ و الی آخر؛
در واقع، اختلاف اندازه در هر گام فقط \emph{بزرگ‌تر} می‌شود. پس وسوسه
می‌شویم حدس بزنیم که برای هر عدد ترتیبی $\kappa$،
$\kappa< \aleph_\kappa$.

اما در $\ZFC$ این حدس \emph{نادرست} است. در واقع می‌توانیم اثبات کنیم
که \emph{نقاط ثابتِ $\aleph$} وجود دارند، یعنی کاردینال‌هایی چون
$\kappa$ که $\kappa=\aleph_\kappa$.

\begin{prop}\ollabel{alephfixed}
یک نقطهٔ ثابتِ $\aleph$ وجود دارد.
\end{prop}

\begin{proof}
با استفاده از بازگشت تعریف کنید:
\begin{align*}
	\kappa_0 &= 0\\
	\kappa_{n+1} &= \aleph_{\kappa_n}\\
	\kappa&= \bigcup_{n < \omega}\kappa_n
\end{align*}
اکنون بنا بر \olref[cardinals][classing]{unioncardinalscardinal}،
$\kappa$ یک کاردینال است. همچنین:
\[
	\kappa= \bigcup_{n < \omega} \kappa_{n+1} =
	\bigcup_{n < \omega}\aleph_{\kappa_n} =
	\bigcup_{\alpha < \kappa}\aleph_\alpha = \aleph_\kappa
\]
\end{proof}

بولوس زمانی مقاله‌ای دقیقاً دربارهٔ همان نقطهٔ ثابتِ $\aleph$ نوشت که
تازه ساختیم. او پس از اشاره به وجود $\kappa$ در آغاز مقاله‌اش گفت:
\begin{quote}
	[$\kappa$] به معیار کسانی که پیش‌تر هیچ آشنایی‌ای با نظریهٔ
	مجموعه‌ها نداشته‌اند، عددی \emph{بسیار بزرگ} است؛ به نظر من، آن‌قدر
	بزرگ که صدق هر نظریه‌ای را زیر سؤال می‌برد که یکی از گزاره‌هایش مدعی
	وجود دست‌کم $\kappa$ شیء باشد.
	\citep[p.~257]{Boolos2000}
\end{quote}
و سرانجام مقاله را با این پرسش به پایان برد:
\begin{quote}
	[آیا] گمان می‌کنیم که، هرچه در آغاز داستان رخ داده باشد، تا وقتی به
	این‌جا رسیده‌ایم چرخ‌ها بی‌حاصل می‌چرخند و دیگر به وصف هیچ امرِ واقع
	گوش نمی‌دهیم؟
	\citep[p.~268]{Boolos2000}
\end{quote}
اگر واقعاً از «هر آنچه واقع است» فراتر رفته‌ایم، باید انگشت اتهام را
مستقیماً به سوی اصل جایگزینی بگیریم، زیرا همین اصل است که اثبات ما را
ممکن می‌کند. در این صورت، احتمالاً بولوس باید در ادعایی تجدیدنظر کند که
چند دهه پیش‌تر مطرح کرده بود، مبنی بر این‌که جایگزینی هیچ پیامد
«نامطلوبی» ندارد (بنگرید به \olref[replacement][extrinsic]{sec}).

اما آیا وجود $\kappa$ تا این اندازه بد است؟ در این‌جا شاید توجه به
\emph{پارادوکس تریسترام شندی} راسل مفید باشد. تریسترام شندی زندگی خود را
در دفتر خاطراتش ثبت می‌کند، اما نوشتن شرح یک روز برایش یک سال زمان
می‌برد. با گذشت هر سال، تریسترام بیشتر و بیشتر عقب می‌افتد: پس از یک
سال، تنها یک روز را ثبت کرده و ۳۶۴ روز ثبت‌نشده زندگی کرده است؛ پس از دو
سال، تنها دو روز را ثبت کرده و ۷۲۸ روز ثبت‌نشده زندگی کرده است؛ پس از سه
سال، تنها سه روز را ثبت کرده و ۱۰۹۲ روز ثبت‌نشده زندگی کرده است
\dots\footnote{سال‌های کبیسه را نادیده می‌گیریم.} بااین‌همه، اگر
تریسترام \emph{نامیرا} باشد، موفق می‌شود هر روز را ثبت کند، زیرا روز
$n$ام را در سال $n$ام زندگی‌اش خواهد نوشت. پس «در پایان زمان»، دفتر
خاطرات تریسترام کامل خواهد بود.

حال چرا این تصویر باید چندان متفاوت باشد از این اندیشه که $\alpha$ از
$\aleph_\alpha$ کوچک‌تر است---و در واقع به‌طور فزاینده و نومیدکننده‌ای
کوچک‌تر است---تا این‌که به $\kappa$ می‌رسیم و در آن نقطه جبران می‌کنیم
و $\kappa = \aleph_\kappa$؟

این نکته را کنار می‌گذاریم و، با فرض پذیرش $\ZFC$، بحث را با چند نتیجهٔ
سرگرم‌کنندهٔ دیگر دربارهٔ ساختن نقاط ثابت به پایان می‌بریم. سه نتیجهٔ
بعدی، به بیان شهودی، نشان می‌دهند نقطه‌ای غیربدیهی وجود دارد که در آن
پهنای سلسله‌مراتب با بلندی آن برابر است:

\begin{prop}\ollabel{bethfixed}
یک نقطهٔ ثابتِ $\beth$ وجود دارد؛ یعنی $\kappa$ای هست که
$\kappa= \beth_\kappa$.
\end{prop}

\begin{proof}
همانند \olref{alephfixed}، با جایگزین‌کردن «$\beth$» به جای
«$\aleph$».
\end{proof}

\begin{prop}\ollabel{stagesize}
$\card{V_{\omega+\alpha}} = \beth_{\alpha}$. اگر
$\omega \ordtimes \omega \leq \alpha$، آنگاه
$\card{V_\alpha} = \beth_\alpha$.
\end{prop}

\begin{proof}
ادعای نخست با یک استقرای ترامتناهی ساده برقرار است. ادعای دوم از این
واقعیت نتیجه می‌شود که اگر $\omega \ordtimes \omega \leq \alpha$،
آنگاه $\omega + \alpha = \alpha$. برای اثبات این مطلب از واقعیت‌های
حساب ترتیبی در \olref[ord-arithmetic][]{chap} استفاده می‌کنیم. نخست
توجه کنید که
$\omega \ordtimes \omega = \omega \ordtimes (1 \ordplus \omega) =
(\omega \ordtimes 1) \ordplus (\omega\ordtimes\omega) = \omega
\ordplus (\omega \ordtimes \omega)$. اکنون اگر
$\omega \ordtimes \omega \leq \alpha$، یعنی برای یک $\beta$ داشته باشیم
$\alpha = (\omega\ordtimes\omega) \ordplus \beta$، آنگاه
$\omega \ordplus \alpha = \omega \ordplus
((\omega \ordtimes \omega) \ordplus \beta) =
(\omega \ordplus (\omega \ordtimes \omega)) \ordplus \beta =
(\omega \ordtimes \omega) \ordplus \beta = \alpha$.
\end{proof}

\begin{cor}
$\kappa$ای وجود دارد که $\card{V_\kappa} = \kappa$.
\end{cor}

\begin{proof}
بگذارید $\kappa$ نقطهٔ ثابتِ $\beth$ باشد که
\olref{bethfixed} وجودش را تضمین می‌کند. آشکارا
$\omega \ordtimes \omega < \kappa$. پس بنا بر \olref{stagesize}،
$\card{V_\kappa} = \beth_\kappa= \kappa$.
\end{proof}

تعداد مراحلِ زیر $V_\kappa$ با تعداد !!{element}های $V_\kappa$ برابر
است. پس، به‌طور شهودی، پهنای $V_\kappa$ با بلندی آن برابر است. این تصویر
بسیار شبیه تریسترام شندی است: از هر مرحله با گرفتن \emph{مجموعهٔ توانی}
به مرحلهٔ بعد می‌رویم و بدین‌ترتیب سلسله‌مراتب خود را در هر گام
\emph{بسیار} بزرگ‌تر می‌کنیم. اما «در پایان»، یعنی در مرحلهٔ $\kappa$،
پهنای سلسله‌مراتب به بلندی آن می‌رسد.

ممکن است بپرسیم: \emph{پهنای سلسله‌مراتب هر چند وقت یک‌بار با بلندی آن
برابر می‌شود؟} پاسخ این است: \emph{به تعداد اعداد ترتیبی.} اما این پاسخ
به اندکی توضیح نیاز دارد.

عبارت $\tau$ را چنین تعریف می‌کنیم. برای هر $A$، بگذارید:
\begin{align*}
	\tau_0(A) & \defis \cardsucc{\card{A}}\\
	\tau_{n+1}(A) & \defis \beth_{\tau_n(A)}\\
	\tau(A) & \defis \bigcup_{n < \omega}\tau_n(A)
\intertext{همانند \olref{bethfixed}، برای هر $A$، $\tau(A)$ یک نقطهٔ
ثابتِ $\beth$ است و به‌سبب تعریف $\tau_0(A)$ آشکارا
$\card{A} < \tau(A)$. اکنون تعریف بازگشتی زیر را در نظر بگیرید:}
	W_0 &\defis 0\\
	W_{\alpha + 1} & \defis \tau(W_\alpha)\\
	W_\alpha & \defis \bigcup_{\beta < \alpha} W_\beta
	\text{، هرگاه $\alpha$ حدی باشد.}
\end{align*}
این ساخت برای همهٔ اعداد ترتیبی تعریف شده است. چون در هر گام جانشین
$\card{W_\alpha}<\tau(W_\alpha)=W_{\alpha+1}$ و در مراحل حدی اجتماع
گرفته می‌شود، نگاشت $\alpha\mapsto W_\alpha$ اکیداً صعودی است؛ پس، به
بیان شهودی، $W$ یک «!!a{injection}» از اعداد ترتیبی به نقاط ثابتِ
$\beth$ است. و درست همانند قبل، برای هر $\alpha$، پهنای
$V_{W_\alpha}$ با بلندی آن برابر است.

\end{document}
